\subsubsection{Overview}
\label{sec:nf-overview}
This overview is adapted from \hlhref{https://en.wikipedia.org/wiki/Database_normalization\#Normal_forms}{this Wikipedia article}.
All constraint definitions in {\scriptsize smaller font} are informal. Note that \textit{prime attributes} are the attributes that make up a candidate key,
also known as \textit{key attributes}

\newcommand{\cross}{\ding{55}}
\begin{tables}{p{10cm}lllll}{Constraint & UNF        & 1NF        & 2NF        & 3NF        & BCNF}
              Unique rows {\scriptsize (no duplicates)}
                                        & \checkmark & \checkmark & \checkmark & \checkmark & \checkmark \\
              Scalar Columns {\scriptsize (Columns can't contain relations or composite values)}
                                        & \cross     & \checkmark & \checkmark & \checkmark & \checkmark \\
              Every non-prime attribute has a full FD on each candidate key {\scriptsize (Attributes depend on the whole of every key)}
                                        & \cross     & \cross     & \checkmark & \checkmark & \checkmark \\
              Every Non-trivial FD begins either with superkey or ends with prime attribute {\scriptsize (Attributes depend only on candidate keys)}
                                        & \cross     & \cross     & \cross     & \checkmark & \checkmark \\
              Every non-trivial FD begins with a superkey
                                        & \cross     & \cross     & \cross     & \cross     & \checkmark \\
\end{tables}
This means that a relation in BNCF is free of redundancies based on functional dependencies.

Note that the above table needs to be read vertically, i.e. for UNF only Unique Rows needs to be satisfied,
thus it doesn't mean that if we have Scalar Columns that this isn't allowed in UNF. In fact, it is allowed, as it is a more strict property.
